%% ----------------------------------------------------------------------------
%%  §4.9  Appendice: dimostrazione del teorema di Weierstrass
%% ----------------------------------------------------------------------------
\section{Appendice: dimostrazione del teorema di Weierstrass}
\label{sec:dim-weierstrass}

Dimostriamo ora il teorema enunciato nel \cref{sec:teoremi-continue}: \emph{se $f$ è una
funzione continua in un intervallo chiuso e limitato $[a,b]$, allora $f$ assume minimo e
massimo in $[a,b]$}, cioè esistono $x_{1},x_{2}\in[a,b]$ per cui vale la
\eqref{eq:weierstrass}.

\begin{strategia}[Piano d'attacco per il \cref{teo:weierstrass}]
Consideriamo l'estremo superiore $M$ dei valori di $f$, che a priori potrebbe anche
essere $+\infty$. Per le proprietà dell'estremo superiore esiste una successione di punti
$x_{n}$ in cui i valori $f(x_{n})$ si avvicinano a $M$ (una successione
\emph{massimizzante}). I punti $x_{n}$ potrebbero non convergere, ma stanno tutti in
$[a,b]$, che è limitato: per il teorema di Bolzano--Weierstrass se ne può estrarre una
successione convergente a un punto $x_{0}$, che sta ancora in $[a,b]$ perché
l'intervallo è chiuso. Per la continuità, i valori di $f$ lungo l'estratta tendono a
$f(x_{0})$; ma tendono anche a $M$. Dunque $M=f(x_{0})$: l'estremo superiore è finito ed
è un valore assunto, cioè un massimo. Ogni ipotesi interviene in un punto preciso: la
limitatezza dell'intervallo per applicare Bolzano--Weierstrass, la chiusura per non
perdere il punto limite, la continuità per passare al limite dentro $f$.
\end{strategia}

\begin{dimostrazione}[Dimostrazione del \cref{teo:weierstrass}]
Posto $M=\sup\{f(x) : x\in[a,b]\}$ (con la convenzione $M=+\infty$ se l'insieme dei
valori non è limitato superiormente, \eqref{eq:sup-infinito}), verifichiamo che esiste
una successione $x_{n}$ di punti di $[a,b]$ tale che
\begin{equation}\label{eq:weier-massimizzante}
  \lim_{n\to+\infty}f(x_{n})=M .
\end{equation}
Infatti, se $M=+\infty$, nessun numero $n\in\N$ è un maggiorante dei valori di $f$, e
quindi per ogni $n\in\N$ esiste $x_{n}\in[a,b]$ tale che $f(x_{n})>n$; perciò
$f(x_{n})\to M=+\infty$ (\cref{teo:confronto-infinito}). Se invece risulta $M<+\infty$,
per la caratterizzazione dell'estremo superiore (\cref{prop:caratt-sup}) con
$\epsilon=\frac1n$, per ogni $n\in\N$ esiste $x_{n}$ in $[a,b]$ tale che
\begin{equation}\label{eq:weier-stima}
  M-\frac1n<f(x_{n})\le M,
\end{equation}
e perciò $f(x_{n})\to M$ per il teorema dei carabinieri.

La successione $x_{n}$ è limitata, perché $a\le x_{n}\le b$ per ogni $n$. Per il teorema
di Bolzano--Weierstrass (\cref{teo:bolzano-weierstrass}) esistono un'estratta $x_{n_{k}}$
da $x_{n}$ e un punto $x_{0}$ tali che
\begin{equation}\label{eq:weier-estratta}
  x_{n_{k}}\to x_{0} \qquad\text{per } k\to+\infty ;
\end{equation}
inoltre $x_{0}\in[a,b]$, per il \cref{cor:confronto-limiti} applicato alle disuguaglianze
$a\le x_{n_{k}}\le b$. Poiché $f$ è continua in $x_{0}$, ne segue
(\cref{prop:caratt-continuita}(b), che ammette anche eventuali termini $x_{n_{k}}=x_{0}$)
\begin{equation}\label{eq:weier-continuita}
  f(x_{n_{k}})\to f(x_{0}) .
\end{equation}
D'altra parte $f(x_{n_{k}})$ è un'estratta di $f(x_{n})$, e quindi, per la
\eqref{eq:weier-massimizzante} e la \cref{prop:estratta-limite}, tende a $M$. Per
l'unicità del limite
\begin{equation}\label{eq:weier-uguaglianza}
  M=\lim_{n\to+\infty}f(x_{n})=\lim_{k\to+\infty}f(x_{n_{k}})=f(x_{0}) .
\end{equation}
Abbiamo così dimostrato che
\begin{equation}\label{eq:weier-tesi}
  f(x_{0})=M=\sup\{f(x) : x\in[a,b]\};
\end{equation}
ciò implica allo stesso tempo che $M<+\infty$ (perché $f(x_{0})$ è un numero reale) e
che l'estremo superiore è, in effetti, un massimo: posto $x_{2}=x_{0}$, si ha
$f(x)\le f(x_{2})$ per ogni $x\in[a,b]$.

Analogamente si ragiona per determinare un punto di minimo, partendo dall'estremo
inferiore di $f$ in $[a,b]$; oppure si applica quanto appena dimostrato alla funzione
continua $-f$: un punto di massimo $x_{1}$ di $-f$ verifica $-f(x)\le-f(x_{1})$, cioè
$f(x_{1})\le f(x)$, per ogni $x\in[a,b]$.
\end{dimostrazione}

\begin{corollario}\label{cor:continua-limitata}
Una funzione continua in un intervallo chiuso e limitato è limitata.
\end{corollario}

\begin{schema}
\Ipotesi{$f$ continua in $[a,b]$.}
\Tesi{$\esiste L>0$: \ $\abs{f(x)}\le L$ \ $\perogni x\in[a,b]$.}
\end{schema}

\begin{strategia}
Il teorema di Weierstrass dà molto di più della limitatezza: un minimo e un massimo
\emph{reali}. Basta allora racchiudere l'intervallo $[m,M]$ in un intervallo simmetrico
$[-L,L]$.
\end{strategia}

\begin{dimostrazione}
Per il teorema di Weierstrass $m\le f(x)\le M$ per ogni $x\in[a,b]$, con $m$, $M$ numeri
reali; quindi $\abs{f(x)}\le L=\max\{\abs m,\abs M,1\}$ per ogni $x\in[a,b]$
(\cref{prop:limitato-modulo}).
\end{dimostrazione}

Rileggiamo alla luce della dimostrazione gli esempi del \cref{sec:teoremi-continue}
(\cref{fig:controesempi-weierstrass}). Per $\frac1x$ su $(0,1]$ e per $g$ su $[-1,1]$
l'estremo da raggiungere (rispettivamente $+\infty$ e $0$) è avvicinato lungo una
successione che tende a $0$: nel primo caso il punto limite $0$ esce dall'intervallo, che
non è chiuso; nel secondo il punto limite c'è, ma la discontinuità di $g$ impedisce di
passare al limite dentro la funzione, e $g(0)=1$ non ha nulla a che fare con $0$. Per
$\frac1x$ su $[1,+\infty)$ la successione minimizzante $x_{n}=n$ diverge, e nessuna sua
estratta converge: è la limitatezza dell'intervallo che viene a mancare.

\begin{notastorica}[Bolzano, Weierstrass e la compattezza]
Karl Weierstrass dimostrò il teorema che porta il suo nome nelle lezioni berlinesi dei
primi anni Sessanta dell'Ottocento, con un argomento basato proprio su quello che oggi chiamiamo teorema di
Bolzano--Weierstrass. Ancora una volta, però, Bolzano lo aveva preceduto: nella sua
\emph{Functionenlehre}, scritta intorno al 1830, si trova una dimostrazione del fatto che
una funzione continua su un intervallo chiuso e limitato è limitata e assume il massimo;
il manoscritto rimase inedito fino al 1930. La proprietà dell'intervallo $[a,b]$ usata
nella dimostrazione, che \emph{ogni successione di suoi punti ammette un'estratta
convergente a un punto dell'intervallo}, fu isolata all'inizio del Novecento da Maurice
Fréchet, che tra il 1904 e il 1906 introdusse per gli insiemi con questa proprietà il
nome di \emph{compatti} (nella sua terminologia originaria il punto limite non doveva
necessariamente appartenere all'insieme). Con questo linguaggio il teorema di Weierstrass si enuncia in grande
generalità: una funzione continua su un insieme compatto ha massimo e minimo.
\end{notastorica}
